% ============================================================
%  capitulo4_tecnico.tex
%  Capítulo 4: Análisis Técnico y Financiero de la
%              Implementación de AGPE en Normalización
% ============================================================

\chapter{Análisis Técnico y Financiero de la Implementación de AGPE en Normalización}
\label{chap:tecnico}

Este capítulo tiene dos partes. La primera revisa el comportamiento
histórico de la normalización en Colombia desde la perspectiva
financiera, cruzando los datos de inversión privada de Air-e y Afinia
con el registro de proyectos del PRONE entre 2005 y 2025. La segunda
desarrolla el caso de estudio hipotético que evalúa la viabilidad
técnica y financiera de integrar sistemas AGPE en un proyecto de
normalización bajo las condiciones reales de la región Caribe.

Los parámetros de demanda y las condiciones climáticas y
socioeconómicas del caso de estudio vienen del
Capítulo~\ref{chap:diagnostico}. El marco normativo de referencia
está en el Capítulo~\ref{chap:normativo}.


% ─────────────────────────────────────────────────────────────
\section{Normalización a Través de los Años con Recursos Públicos y Privados}
\label{sec:historico}
% ─────────────────────────────────────────────────────────────

La literatura sobre normalización en Colombia documenta poco lo que
pasa en campo. No hay muchos estudios que registren un proyecto desde
la pre-factibilidad hasta la operación con datos suficientes para
analizarlo. Lo que sí hay son los casos del Capítulo~\ref{chap:normativo},
las actas CAPRONE del MME y los informes quinquenales de Air-e y Afinia
para los años 2021--2024. Con eso se construyó este análisis.

\subsection{Comportamiento de la normalización con recursos privados}
\label{subsec:rec_privados}

Los dos casos con inversión privada documentada son los de Santander
(2010--2014) y Cauca (2014--2018), analizados en el
Capítulo~\ref{chap:normativo}. Fuera de esos, hay dos estudios con
datos cuantitativos relevantes.

\textcite{CabralesSanjuan2015} estimaron un costo promedio de
\textbf{\$1.093.052 por usuario} en proyectos de normalización urbana
en CENS, orientados principalmente a la sustitución de medidores y
acometidas obsoletas. Es la referencia numérica más citada en la
literatura para proyectos de este tipo. \textcite{PenaMantilla2020},
por su parte, analizaron alternativas de mitigación de pérdidas en EPM
para el municipio de Itagüí y encontraron que legalizar y controlar
usuarios es más rentable que normalizar completamente la infraestructura,
aunque advierten que sin revisiones y controles posteriores los usuarios
normalizados tienden a reincidir en las conexiones ilegales.

Al revisar los planes quinquenales de Air-e y Afinia, lo que se ve es
una priorización constante de los niveles de tensión 2, 3 y 4. El nivel
1 —que es donde está la infraestructura de los BS— queda relegado.

\subsubsection{Air-e: comportamiento de la inversión 2021--2024}
\label{subsubsec:aire_inversion}

La inversión de Air-e en nivel de tensión 1 fue inferior al 30\% durante
2021--2023: 13\% en 2021, 24\% en 2022 y 27\% en 2023. En 2024 subió al
50\%, pero se destinó a modernizar líneas y circuitos, no a normalizar
BS. Los informes no identifican ninguna partida específica para usuarios
subnormales \parencite{Aire2021, Aire2022, Aire2023, Aire2024}.

Más preocupante aún: la ejecución del plan de inversiones aprobado por
la CREG fue del 25\% en 2023 y del 19\% en 2024. El OR justificó esos
incumplimientos en dificultades operativas y problemas de liquidez. La
conclusión es directa: si Air-e no puede ejecutar ni su propio plan de
inversiones, su capacidad de cofinanciar proyectos de normalización con
recursos propios es prácticamente nula. El PRONE no es un complemento
conveniente para ese OR — es la única vía realista.

\subsubsection{Afinia: estrategias de control de pérdidas 2021--2024}
\label{subsubsec:afinia_perdidas}

Afinia no detalla en sus informes la inversión por nivel de tensión, pero
sí muestra un esfuerzo sostenido en reducción y control de pérdidas
mediante acciones técnicas y comerciales: normalización de medida en
suministros con altas pérdidas, medición prepago, adecuación de redes
internas y talleres de gestión social para mejorar el recaudo
\parencite{Afinia2021, Afinia2022, Afinia2023, Afinia2024}.

La Tabla~\ref{tab:afinia_acciones} muestra las acciones ejecutadas entre
2021 y 2024. Las pérdidas bajaron de 27,69\% en 2021 a 26,15\% en 2022,
pero volvieron a subir en 2023 (26,73\%) y 2024 (27,95\%). El OR
atribuye el rebote al aumento del hurto de energía por la inconformidad
de los usuarios con el alza del Costo Unitario (CU) por encima de
\$1.000/kWh, y al mayor consumo generado por el Fenómeno de El Niño.

\begin{table}[H]
\centering
\caption{Acciones realizadas por Afinia para el control y reducción de pérdidas (2021--2024)}
\label{tab:afinia_acciones}
\begin{tabularx}{\textwidth}{>{\RaggedRight}X
                             >{\centering}p{2cm}
                             >{\centering}p{2cm}
                             >{\centering}p{2cm}
                             >{\centering\arraybackslash}p{2cm}}
\toprule
\textbf{Acción} & \textbf{2021} & \textbf{2022} & \textbf{2023} & \textbf{2024} \\
\midrule
\rowcolor{grisclaro}
Normalización de suministros sin medidor            & 16.841 & 10.580 & 13.260 & 8.487  \\
Suministros con medición inteligente                & 13.738 & 26.723 & ---    & ---    \\
\rowcolor{grisclaro}
Instalación de macromedidores NT1                   & 281    & 740    & 2.739  & 10.851 \\
Incorporación de nuevas invasiones (contratación)   & 86     & 34     & ---    & ---    \\
\rowcolor{grisclaro}
Incorporación de clientes conectados ilícitamente   & 6.145  & 36.866 & 35.118 & 36.774 \\
Normalización de la medida                          & 1.719  & 6.250  & 14.131 & 19.461 \\
\rowcolor{grisclaro}
Instalación de equipos de medida prepago            & 255    & 270    & ---    & ---    \\
\midrule
\textbf{Total}  & \textbf{39.065} & \textbf{81.463} & \textbf{65.248} & \textbf{75.573} \\
\bottomrule
\multicolumn{5}{l}{\footnotesize\textit{Fuente: Elaboración propia
con datos de \parencite{Afinia2021, Afinia2022, Afinia2023, Afinia2024}.}}
\end{tabularx}
\end{table}

La instalación de macromedidores en NT1 pasó de 281 en 2021 a 10.851 en
2024, y la normalización de la medida de 1.719 a 19.461. Hay aprendizaje
y hay escala creciente. Pero el rebote en pérdidas de los últimos dos
años muestra que las estrategias implementadas no son suficientes para
sostener la tendencia. Las limitaciones operativas y financieras siguen
presentes.

\subsection{Comportamiento de la normalización con recursos públicos}
\label{subsec:rec_publicos}

La base de datos construida a partir de las actas CAPRONE del MME
registra, por vigencia, el número de proyectos presentados por OR, el
valor aprobado y los usuarios beneficiarios \parencite{MME2025}. El
archivo complementario \textit{Proyectos adjudicados PRONE 2005--2025}
contiene el detalle completo por municipio y departamento.

Con esa información se construyó el análisis en Power BI que se presenta
en las figuras siguientes.

\begin{figure}[H]
\centering
\includegraphics[width=\textwidth]{figura_prone_historico.png}
\caption{Proyectos financiados con recursos públicos, vigencias 2003--2025}
\label{fig:prone_historico}
\footnotesize\textit{Fuente: Elaboración propia con datos de \parencite{MME2025}.}
\end{figure}

La concentración geográfica es clara: región Caribe, región Pacífica y
algunos departamentos andinos dominan el mapa. El 52,03\% de los
proyectos entre 2005 y 2025 provienen de la Electrificadora de la Costa
Caribe (hoy Afinia), seguida por Air-e (Figura~\ref{fig:prone_or}).

\begin{figure}[H]
\centering
\includegraphics[width=0.72\textwidth]{figura_prone_por_or.png}
\caption{Proyectos presentados por Operador de Red, vigencias 2005--2025}
\label{fig:prone_or}
\footnotesize\textit{Fuente: Elaboración propia con datos de \parencite{MME2025}.}
\end{figure}

Esa concentración no garantiza calidad en la formulación. En 2024,
ningún proyecto de ningún OR cumplió los requisitos. El caso de CENS
ilustra el problema desde otro ángulo: ese OR no ha participado en las
convocatorias del MME porque las alcaldías de su área de cobertura no
han ejecutado las certificaciones de BS requeridas
\parencite{CabralesSanjuan2015}. La barrera no es técnica. Es
institucional.

\subsubsection{Variaciones anuales en proyectos y beneficiarios}

La Figura~\ref{fig:prone_proyectos_usuarios} muestra el número de
proyectos y usuarios beneficiados por vigencia.

\begin{figure}[H]
\centering
\includegraphics[width=0.90\textwidth]{figura_prone_proyectos_usuarios.png}
\caption{Proyectos presentados y usuarios beneficiados, vigencias 2005--2025}
\label{fig:prone_proyectos_usuarios}
\footnotesize\textit{Fuente: Elaboración propia con datos de \parencite{MME2025}.}
\end{figure}

Entre los periodos más relevantes se destacan los siguientes:

\textbf{2007--2010}: Sin proyectos aprobados. El ajuste normativo tras
la Ley 1117 de 2006 y el Decreto 1123 de 2008 detuvo el programa
mientras se redefinían condiciones, actores y procedimientos.

\textbf{2012}: 185 proyectos, el pico histórico del programa. Resultado
probable de años de acumulación de diagnósticos técnicos en los OR.

\textbf{2016--2018}: De nuevo sin adjudicaciones. Coincide con el ajuste
institucional del DUR 1073 de 2015.

\textbf{2024}: Cero proyectos elegibles. Primera convocatoria sin
ningún resultado. Ese dato confirma la brecha metodológica que justifica
esta monografía \parencite{MME2025}.

\textbf{2025}: 30 proyectos con el nuevo esquema AGPE, beneficiando
cerca de 11.000 usuarios en La Guajira, Magdalena, Atlántico, Nariño,
Huila y Meta.

\subsubsection{Inversión histórica por vigencia}

\begin{figure}[H]
\centering
\includegraphics[width=0.65\textwidth]{figura_prone_inversion.png}
\caption{Inversión por año en proyectos de normalización, vigencias 2005--2025}
\label{fig:prone_inversion}
\footnotesize\textit{Fuente: Elaboración propia con datos de \parencite{MME2025}.}
\end{figure}

\begin{figure}[H]
\centering
\includegraphics[width=\textwidth]{figura_prone_2025.png}
\caption{Proyectos financiados en la vigencia 2025}
\label{fig:prone_2025}
\footnotesize\textit{Fuente: Elaboración propia con datos de \parencite{MME2025}.}
\end{figure}

La vigencia 2025 es la más atípica del periodo analizado: con solo 30
proyectos —menos de una quinta parte del pico de 2012— el valor de
financiación solicitado es comparable al de vigencias con muchos más
proyectos. Eso es directamente atribuible al componente AGPE: cada
sistema fotovoltaico representa un costo adicional que eleva de manera
sustancial el costo por usuario respecto a la normalización convencional.

En total, entre 2003 y 2025 el PRONE reporta cerca de \textbf{290.000
usuarios beneficiados} según proyectos aprobados. Esa cifra puede
diferir del número real ejecutado por los ajustes que ocurren durante
la obra.


% ─────────────────────────────────────────────────────────────
\section{Análisis Técnico del Caso de Estudio}
\label{sec:caso_tecnico}
% ─────────────────────────────────────────────────────────────

Esta sección construye un caso de estudio hipotético basado en las
condiciones reales de la región Caribe para evaluar la viabilidad de
integrar AGPE en un proyecto de normalización. Los parámetros de
demanda vienen del Capítulo~\ref{chap:diagnostico} y la normativa de
referencia del Capítulo~\ref{chap:normativo}. El análisis compara los
escenarios con y sin AGPE.

\subsection{Lineamientos técnicos y sociales para la integración de sistemas AGPE en proyectos de normalización}
\label{subsec:lineamientos_agpe}

La integración de estos dos sistemas eléctricos que garanticen un beneficio real y duradero para los usuarios necesita el cumplimiento de estándares técnicos y el fortalecimiento comunitario y social en pro de la apropiación y sentido de pertenencia sobre el proyecto. A continuación se definen los lineamientos mínimos para ello.

\subsubsection{Lineamientos Técnicos}

\textbf{Articulación normativa:} El Reglamento Técnico de Instalaciones Eléctricas (RETIE) establece las condiciones mínimas que una instalación eléctrica debe cumplir, entre ellas distancias de seguridad, selección de constructores y regulación de tensión, entre otros; y la Resolución CREG 174 de 2021 constituyen la normatividad que rige el diseño, dimensionamiento y la conexión del sistema AGPE a la red de distribución. Dentro de los requisitos principales para el dimensionamiento, la Resolución 174 limita la potencia máxima instalada de un sistema AGPE al 50\% de la capacidad del transformador o circuito de conexión \parencite{ResolucionCREG174_2021}.

Adicionalmente, estipula que los excedentes que superen la importación se reconocerán a los usuarios a través de créditos de energía. Aunado a ello, establece los componentes de la tarifa que el usuario debe pagar según la capacidad del sistema AGPE ($<$ ó $>$ 100~kW): Comercialización, Distribución, Transmisión, Pérdidas Reconocidas y Restricciones. Así mismo, es pertinente recalcar que si bien ambos sistemas están integrados, el dimensionamiento de la red de distribución no está estrictamente ligado al del sistema AGPE, dado que si el sistema es de generación solar fotovoltaica su mayor producción se da en horas del día, mientras que la curva de carga muestra que el pico máximo de demanda se da en horas de la noche. En consecuencia, el diseño de la red de distribución sigue estando ligado principalmente a los picos de demanda \parencite{CREGCargoRespaldo2023, ResolucionCREG174_2021}.

\textbf{Sistema de medición inteligente:} Debido a la naturaleza del proyecto y a la exigencia en las convocatorias del MME, se requieren medidores con tecnología AMI y bidireccionales, ya que permiten identificar los excedentes entregados a la red por el sistema AGPE con el fin de ser reconocidos a los usuarios.

\textbf{Dimensionamiento del sistema de AGPE:} Con base en las dos resoluciones direccionadas al financiamiento de sistemas de normalización con AGPE, la capacidad instalada de los sistemas debe garantizar el consumo básico de subsistencia. Para ello, el cálculo de potencia debe ajustarse al factor de planta real del área de influencia.

\subsubsection{Lineamientos Sociales}

\begin{itemize}
    \item \textbf{Conformación de Comunidades Energéticas:} A través del ``Plan de Acción Social'' formulado por el MME en las convocatorias de la vigencia 2025, la organización de los usuarios bajo la figura de Comunidades Energéticas es un requisito durante la ejecución de los proyectos. Si bien no se define el prestador de servicio público responsable del AOM del sistema de AGPE, este debe plantearse teniendo en cuenta que la gestión del sistema es colectiva y es importante el fortalecimiento y apropiación de los usuarios con respecto a la infraestructura eléctrica del proyecto.

    \item \textbf{Formación y capacitación:} Conforme al Plan de Acción, este tipo de proyecto hace obligatoria la formación básica de los beneficiarios en AOM, en el uso racional y gestión eficiente de la energía. Esto engloba temas como la interpretación de los datos de un medidor y la cadena de valor. Los conocimientos adquiridos podrán romper las barreras sociales que afectan la ejecución de los proyectos de normalización en torno al aumento de la factura, mejorando los hábitos de consumo en los BS y garantizando la sostenibilidad del sistema de AGPE.
\end{itemize}


\subsection{Dimensionamiento técnico y económico de la normalización con AGPE}
\label{subsec:dimensionamiento_caso}

El dimensionamiento de la red de distribución integrada y del sistema de autogeneración se realiza teniendo en cuenta, entre otras variables, la demanda energética, la capacidad de los equipos de generación fotovoltaica y la infraestructura mínima de una red en Media Tensión (MT) y Baja Tensión (BT) requerida para garantizar el cumplimiento de la normativa vigente.

Para la selección del área de influencia, se tomaron como referencia los proyectos presentados en las convocatorias del año 2025, dado que representan el único antecedente a la fecha de proyectos de normalización adjudicados con inclusión de sistemas AGPE. De los cuatro OR que presentaron proyectos y obtuvieron concepto FAVORABLE, Air-e fue el único con presencia en la región Caribe. Teniendo en cuenta que en el Capítulo~\ref{chap:diagnostico} se determinó que esta región es la zona donde más se presenta el fenómeno de subnormalidad y que el clima cálido húmedo eleva el consumo por necesidades de confort térmico, se seleccionó el proyecto en el BS ``El Espinal''.

A continuación se describen los criterios de entrada:

\begin{itemize}
    \item \textbf{Ubicación:} municipio Malambo, Atlántico
    \item \textbf{Cantidad de usuarios:} 106
    \item \textbf{Parámetros ambientales promedio:} T~27\,°C; HR~82\%; A~10~msnm
    \item \textbf{Consumo promedio por usuario:} 235~kWh/mes
    \item \textbf{Consumo básico de subsistencia:} 173~kWh/mes
    \item \textbf{Estrato socioeconómico:} 1
    \item \textbf{Certificación:} expedida por el ente territorial como BS
    \item \textbf{Área disponible:} El polígono resaltado en rojo tiene un área de 77,72~ha, que para el caso de estudio corresponde al BS ``El Espinal''. Así mismo, la ubicación de las 106 viviendas está alrededor de la ruta marcada en amarillo, la cual tiene una longitud de 1,289~km (ver Figura~\ref{fig:delimitacion_bs}).
\end{itemize}

\begin{figure}[H]
\centering
%\includegraphics[width=0.85\textwidth]{figura_delimitacion_bs.png}
\caption{Delimitación geográfica del BS ``El Espinal''}
\label{fig:delimitacion_bs}
\footnotesize\textit{Fuente: Elaboración propia en Google Earth.}
\end{figure}


\subsection{Diseño eléctrico de la normalización de redes}
\label{subsec:disenio_electrico}

El diseño eléctrico del proyecto se desarrolla conforme a la Resolución
MME 40117 de 2024 que modificó el RETIE. La normatividad técnica de
referencia es:

\begin{itemize}
    \item RETIE -- Resolución MME 40117 de 2024.
    \item RETILAP.
    \item Normas ICONTEC.
    \item Código Eléctrico Colombiano (NTC 2050).
    \item Normas y especificaciones del OR correspondiente.
    \item Especificaciones de fabricantes de equipos y materiales.
\end{itemize}

Los ítems del diseño, conforme al RETIE, son:

\begin{enumerate}[label=\alph*), leftmargin=1.8cm, itemsep=3pt]
    \item Análisis de riesgos de origen eléctrico y medidas de mitigación.
    \item Análisis de riesgos por descargas atmosféricas y medidas de protección.
    \item Análisis y cálculo de cargas iniciales y futuras, incluyendo factor de potencia y armónicos.
    \item Coordinación de aislamiento eléctrico.
    \item Análisis y cálculo de cortocircuito, arco eléctrico y falla a tierra.
    \item Análisis del nivel de tensión requerido.
    \item Cálculos de campos electromagnéticos.
    \item Cálculo de transformadores incluyendo efectos de armónicos y factor de potencia.
    \item Sistema de puesta a tierra.
    \item Cálculo económico de conductores teniendo en cuenta factores de pérdidas y costos de la energía.
    \item Especificación de conductores conforme a IEC 60909 u equivalente.
    \item Cálculo mecánico de estructuras y elementos de sujeción de redes.
    \item Cálculo y coordinación de protecciones contra sobrecorrientes conforme a IEC 60947-2 Anexo A.
    \item Cálculos de canalizaciones, bandejas portacables y volumen de encerramientos.
    \item Cálculo de pérdidas de energía incluyendo efectos de armónicos y factor de potencia.
    \item Cálculos de regulación de tensión.
    \item Identificación de áreas clasificadas como peligrosas.
    \item Diagramas unifilares.
    \item Planos eléctricos para construcción.
    \item Especificaciones técnicas de equipos y materiales.
    \item Distancias de seguridad o servidumbre requeridas.
    \item Justificación de desviaciones técnicas cuando sea necesario.
    \item Selección, cálculo y especificación de equipos de generación convencionales y no convencionales.
\end{enumerate}


\subsection{Dimensionamiento eléctrico y económico del sistema de AGPE}
\label{subsec:dimensionamiento_agpe}

Considerando que el PRONE no especifica el tipo de generación a implementarse, sino por el contrario lo deja abierto a FNCER, fue necesario realizar un análisis de alternativas para identificar la tecnología adecuada para el caso de estudio, teniendo en cuenta el recurso energético disponible, los costos de inversión inicial, los costos fijos de operación y mantenimiento (O\&M) y la generación requerida. El proyecto debe suplir como mínimo el consumo básico de subsistencia (173~kWh/mes, según lo establecido en la Resolución 355 de 2004) de acuerdo con las Resoluciones 40243 y 40432 de 2025; por lo tanto, el sistema se dimensionará con respecto a este valor para los 106 usuarios \parencite{ResolucionCREG355_2004, ResolucionMME40243_2025, ResolucionMME40432_2025}.

Al observar la tabla comparativa de las seis tecnologías dispuestas a pequeña escala (Tabla~\ref{tab:tecnologias_fnce}), se evidencia que en Colombia la tecnología con mayor despliegue es la solar fotovoltaica. Esta opción resulta ser la más económica en términos globales y la que requiere un menor porcentaje de inversión destinado a la instalación, alineado a que requieren un menor tiempo para su construcción, permitiendo que los proyectos se ejecuten con agilidad tal que puedan estar operativos en paralelo a la normalización de las redes eléctricas. Adicionalmente, presenta una ventaja competitiva ya que el costo de O\&M es mucho menor en comparación con las demás tecnologías \parencite{UPMECatalogoTecnologico2024}.

Al analizar las demás tecnologías según sus particularidades, se evidenció que las Pequeñas Centrales Hidroeléctricas (PCH) a filo de agua se ven seriamente afectadas por el fenómeno del Niño debido a la falta de almacenamiento de agua. Por tanto, la disminución del caudal de los ríos durante estos periodos reduciría la capacidad de autogeneración precisamente cuando más se necesita, ya que es en esos instantes cuando la tarifa de energía sube en todo el territorio nacional debido a que la matriz energética depende en más de un 70\% del recurso hídrico \parencite{XMGeneracionReal2025}. Por lo tanto, la autogeneración con PCH no estaría alineada con los picos de tarifa y no atacaría el problema central: la reducción de costos para los usuarios en condiciones de subsistencia durante los meses críticos \parencite{UPMECatalogoTecnologico2024}.

En cuanto a la geotermia, la dinámica de financiación del PRONE exige que en el momento de la convocatoria el sistema ya cuente con un diseño definido. Esto implicaría que el OR debería asumir previamente los gastos de exploración, pruebas y estudios técnicos requeridos para garantizar el recurso. Como se observa en la caracterización técnica y económica, el costo de estos estudios iniciales es elevado, llegando a superar incluso la inversión total de un sistema solar completo, lo que vuelve inviable esta alternativa para un proyecto de normalización social, donde el OR busca financiamiento del Estado para solventar una problemática que afecta su cartera \parencite{UPMECatalogoTecnologico2024}.

De igual manera, la biomasa y el biogás dependen estrictamente de ciclos agrícolas, recolección de residuos y procesos de tratamiento o transporte de insumos. Al tratarse de un municipio con baja actividad económica agrícola, bovina y porcina, especialmente en el área de influencia, no se tendría garantizado el suministro constante de combustible, convirtiéndose en una carga logística y económica para los mismos usuarios \parencite{UPMECatalogoTecnologico2024}.

Finalmente, es importante considerar que todas estas tecnologías alternativas requieren un conocimiento técnico más complejo para su operación y mantenimiento. Dadas las capacidades educativas de los habitantes en estos barrios, un sistema difícil de operar generaría un cuello de botella que afectaría la infraestructura a mediano plazo, especialmente si el sistema queda a cargo de la comunidad bajo la figura de una Comunidad Energética. Por estas razones, la tecnología solar se consolida como la opción más viable, sostenible y de fácil apropiación para los usuarios, aunado a que el municipio de Malambo cuenta con una irradiancia solar en promedio entre 5,0 y 5,5~kWh/m\textsuperscript{2}/día, siendo de los mejores promedios en todo el territorio nacional \parencite{UPMECatalogoTecnologico2024}.

A continuación, en la Tabla~\ref{tab:tecnologias_fnce} se presenta cada una de las tecnologías disponibles en Colombia a pequeña escala y su caracterización técnica y económica.

\begin{table}[H]
\centering
\caption{Comparación técnica y económica de las tecnologías FNCER existentes a pequeña escala}
\label{tab:tecnologias_fnce}
\footnotesize
\begin{tabularx}{\textwidth}{>{\RaggedRight}p{2.2cm}
                             >{\RaggedRight}X
                             >{\RaggedRight}X
                             >{\RaggedRight}X
                             >{\RaggedRight}X
                             >{\RaggedRight}X
                             >{\RaggedRight\arraybackslash}X}
\toprule
\textbf{Atributo} & \textbf{Solar FV} & \textbf{Eólica terrestre} & \textbf{Filo de agua} & \textbf{Geotermia binaria} & \textbf{Biomasa} & \textbf{Biogás} \\
\midrule
\rowcolor{grisclaro}
Categoría & Cat.\ 3, en transición a Cat.\ 4 & Cat.\ 3 & Cat.\ 3 & Cat.\ 3 & Cat.\ 3 & Cat.\ 3 \\

Tamaño y capacidad & 1,5--2,1~m\textsuperscript{2}; 450--670~Wp & 1--25~kW & Hasta 1~MW & 1--10~MW & 1--60~MW & 10~kW--1~MW \\

\rowcolor{grisclaro}
Inversión [MUSD\,2024/MW] & 1,12 (0,23 instalación) & 3,4 (1,0 instalación) & 4,3 (2,86 instalación) & 6,0 (40\% instalación + 2,98 exploración) & 2,56 (35\% instalación) & 4,38 (19\% instalación) \\

O\&M [USD\,2024/MWe] & 11.000/año & 68.500/año & 94.000/año & 88.900/año & 60.700/año & 46.200/año \\

\rowcolor{grisclaro}
Vida útil [años] & 27 & 27 & 50 & 30 & 25 & 25 \\

Construcción [años] & 0,5 & 1 & 2 & 2 & 2 & 1,5 \\

\rowcolor{grisclaro}
Espacio [1000\,m\textsuperscript{2}/MWe] & 12,5 & 11 & 2,4 & 30 & 35 & 0,649 \\

Factor de planta [\%] & 20,3 & 26 & 56 & 88 & 52,6 & 8,2 \\
\bottomrule
\multicolumn{7}{l}{\footnotesize\textit{Fuente: Elaboración propia a partir de \parencite{UPMECatalogoTecnologico2024}.}}
\end{tabularx}
\end{table}

De acuerdo con lo anterior, la demanda de energía diaria a suplir por el sistema fotovoltaico debe ser de mínimo 5,76~kWh/día asumiendo un mes de 30 días. Como se debe suplir mínimo el consumo básico, se toman las condiciones críticas de disponibilidad del recurso solar; por ende, el dimensionamiento se realiza con la irradiancia solar del mes más bajo en el área de influencia.


% ─────────────────────────────────────────────────────────────
\section{Síntesis del Capítulo}
\label{sec:sintesis_cap4}
% ─────────────────────────────────────────────────────────────

El análisis histórico del PRONE entre 2005 y 2025 deja un panorama
claro: en veinte años de operación el programa ha beneficiado a cerca
de 290.000 usuarios, pero la demanda de normalización crece a un ritmo
del 26\% en cuatro años y los OR con mayor exposición a la subnormalidad
no tienen capacidad financiera para cofinanciar proyectos a escala. El
PRONE es la vía prácticamente exclusiva para financiar la normalización
en la región Caribe. El resultado de 2024 —ningún proyecto elegible—
confirma que la brecha no es de recursos sino de capacidad metodológica
para estructurar proyectos que cumplan los requisitos del MME.